
\begin{table}[htbp]
   \caption{\label{tab:main_state} State-Level First-Difference Regressions: Mobilization and Labor Market Outcomes}
   \centering
   \begin{tabular}{lcccccc}
      \tabularnewline \midrule \midrule
      Dependent Variables: & \multicolumn{2}{c}{$\Delta$ Female LFP} & $\Delta$ Female Occ. Score & d\_lf\_gap & $\Delta$ Male LFP & $\Delta$ Female HH Head\\
                               & (1)      & (2)            & (3)      & (4)      & (5)      & (6) \\   
      Model:                   & (1)      & (2)            & (3)      & (4)      & (5)      & (6)\\  
      \midrule
      \emph{Variables}\\
      Mobilization Rate (std.) & -0.0014  & -0.0073$^{**}$ & 0.1481   & -0.0047  & -0.0025  & -0.0011\\   
                               & (0.0038) & (0.0027)       & (0.1329) & (0.0028) & (0.0024) & (0.0015)\\   
      \midrule
      \emph{Fit statistics}\\
      Observations             & 49       & 49             & 49       & 49       & 49       & 49\\  
      R$^2$                    & 0.00285  & 0.67586        & 0.50336  & 0.42320  & 0.47528  & 0.17551\\  
      Adjusted R$^2$           & -0.01837 & 0.62052        & 0.43241  & 0.32472  & 0.40032  & 0.05773\\  
      \midrule \midrule
      \multicolumn{7}{l}{\emph{IID standard-errors in parentheses}}\\
      \multicolumn{7}{l}{\emph{Signif. Codes: ***: 0.01, **: 0.05, *: 0.1}}\\
   \end{tabular}
   
   \par \raggedright 
   Dependent variables are 1940--1950 changes. Column (1): bivariate regression of change in female LFP on standardized mobilization rate. Column (2): adds controls for 1940 state characteristics (percent urban, Black, farm, mean education, mean age, percent married). Column (3): change in female occupation score. Column (4): change in LFP gender gap (female minus male). Column (5): change in male LFP. Column (6): change in female household headship rate. All regressions weighted by 1940 female population (male population for Column 5). N = 48 states plus DC (excluding Alaska and Hawaii).
\end{table}


